\newpage

\gridquestion{13}{En un zoológico, tres animales compiten por ser el más destacado.

La información de cada animal se proporciona mediante un diccionario con cuatro llaves: \texttt{\textquotedbl especie\textquotedbl}, \texttt{\textquotedbl modo\_reproduccion\textquotedbl}, \texttt{\textquotedbl nombre\textquotedbl} y \texttt{\textquotedbl longitud\_patas\textquotedbl}.

E.g.:

\texttt{\{\\
\phantom{xxxx}\textquotedbl especie\textquotedbl: \textquotedbl jirafa\textquotedbl,\\
\phantom{xxxx}\textquotedbl modo\_reproduccion\textquotedbl: \textquotedbl vivíparo\textquotedbl,\\
\phantom{xxxx}\textquotedbl nombre\textquotedbl: \textquotedbl Juanita\textquotedbl,\\
\phantom{xxxx}\textquotedbl longitud\_patas\textquotedbl: 150\\
\}}

Implemente una función que recibe tres animales como parámetros y retorna la especie del animal con las patas más largas. En caso de empate en la longitud de las patas, dé prelación al animal que tenga el nombre propio más largo. En caso de empate total, prefiera el último diccionario en el orden de entrada.

\textit{Consejo: No demore más de 20 minutos en este punto.}}
